% GENERATED FILE. Edit canonical lessons, then run npm run generate:books.
% canonical-source-hash: fnv1a64:5b2e8c62a4b730c7
% canonical-lessons: HI-C404-traifik, HI-C404-jam, HI-C404-motarsaikil, HI-C404-rengna, HI-C404-hakna

\chapter{Traffic Jams}
\label{ch:hi-traffic-jams}
\hlchaptermodality{404}

\begin{chapteropening}
\textbf{By the end of this chapter:} \emph{I can say the traffic is heavy, that we got caught in a jam and that the cars crawled.}

Complete the last lesson of chapter 404: I can say the traffic is heavy, that we got caught in a jam and that the cars crawled.
\end{chapteropening}

\section[\texorpdfstring{\hi{ट्रैफ़िक} (ṭraifik)}{ट्रैफ़िक (ṭraifik)}]{\hi{ट्रैफ़िक} (ṭraifik) — traffic}

\label{lesson:HI-C404-traifik}



\begin{quote}
Before the new one: say the Hindi for to take a stroll, then the Hindi for to pant.
\end{quote}

\subsection*{You'll want to know: \hi{ट्रैफ़िक}}
\textbf{\hi{ट्रैफ़िक}} — \emph{ṭraifik} — \textquotedblleft{}traffic\textquotedblright{}.

\begin{quote}
\hi{शाम} \hi{को} \hi{सड़क} \hi{पर} \hi{बहुत} \hi{ट्रैफ़िक} \hi{होता} \hi{है।}
\end{quote}

\begin{quote}
\emph{śām ko saṛak par bahut ṭraifik hotā hai.}
\end{quote}

\begin{quote}
There is a lot of traffic on the road in the evening.
\end{quote}

Masculine.

\begin{cousinweb}
\textbf{\hi{ट्रैफ़िक}} (\emph{ṭraifik}) is borrowed from English \emph{traffic}, and like its source it takes no plural: \hi{बहुत} \hi{ट्रैफ़िक}, a lot of traffic.
\end{cousinweb}

\begin{grammarlens}[title={What you've built}]
One more word for getting around.
\end{grammarlens}

\subsection*{Your turn}
\begin{itemize}
\raggedright
  \item \emph{Say it:} \emph{ṭraifik}
  \item \emph{Say it:} \emph{ṭraifik}, once more
  \item \emph{Say it:} \emph{śām ko saṛak par bahut ṭraifik hotā hai.} — There is a lot of traffic on the road in the evening.
  \item \emph{Recall:} say the Hindi for to take a stroll, then the Hindi for to pant, then say \emph{ṭraifik} again
\end{itemize}

\begin{tcolorbox}[breakable,colback=teal!4,colframe=teal!35!black,title={Before you move on}]
What does \hi{ट्रैफ़िक} mean? (\textquotedblleft{}Traffic\textquotedblright{}.) Say it once more.
\end{tcolorbox}
\section[\texorpdfstring{\hi{जाम} (jām)}{जाम (jām)}]{\hi{जाम} (jām) — a traffic jam}

\label{lesson:HI-C404-jam}



\begin{quote}
Before the new one: say the Hindi for to pant, then the Hindi for traffic.
\end{quote}

\subsection*{You'll want to know: \hi{जाम}}
\textbf{\hi{जाम}} — \emph{jām} — \textquotedblleft{}a traffic jam\textquotedblright{}.

\begin{quote}
\hi{पुल} \hi{पर} \hi{बड़ा} \hi{जाम} \hi{था।}
\end{quote}

\begin{quote}
\emph{pul par baṛā jām thā.}
\end{quote}

\begin{quote}
There was a big jam on the bridge.
\end{quote}

Masculine.

\hi{जाम} \hi{लगना}, for a jam to build up: \hi{सड़क} \hi{पर} \hi{जाम} \hi{लगा} \hi{है}.

\begin{cousinweb}
\textbf{\hi{जाम}} (\emph{jām}) is borrowed from English \emph{jam}, a road packed so tight that nothing moves.
\end{cousinweb}

\begin{grammarlens}[title={What you've built}]
One more word for getting around.
\end{grammarlens}

\subsection*{Your turn}
\begin{itemize}
\raggedright
  \item \emph{Say it:} \emph{jām}
  \item \emph{Say it:} \emph{jām}, once more
  \item \emph{Say it:} \emph{pul par baṛā jām thā.} — There was a big jam on the bridge.
  \item \emph{Recall:} say the Hindi for to pant, then the Hindi for traffic, then say \emph{jām} again
\end{itemize}

\begin{tcolorbox}[breakable,colback=teal!4,colframe=teal!35!black,title={Before you move on}]
What does \hi{जाम} mean? (\textquotedblleft{}A traffic jam\textquotedblright{}.) Say it once more.
\end{tcolorbox}
\section[\texorpdfstring{\hi{मोटरसाइकिल} (moṭarsāikil)}{मोटरसाइकिल (moṭarsāikil)}]{\hi{मोटरसाइकिल} (moṭarsāikil) — a motorcycle}

\label{lesson:HI-C404-motarsaikil}



\begin{quote}
Before the new one: say the Hindi for traffic, then the Hindi for a traffic jam.
\end{quote}

\subsection*{You'll want to know: \hi{मोटरसाइकिल}}
\textbf{\hi{मोटरसाइकिल}} — \emph{moṭarsāikil} — \textquotedblleft{}a motorcycle\textquotedblright{}.

\begin{quote}
\hi{अरुण} \hi{मोटरसाइकिल} \hi{से} \hi{दफ़्तर} \hi{जाता} \hi{है।}
\end{quote}

\begin{quote}
\emph{Aruṇ moṭarsāikil se daftar jātā hai.}
\end{quote}

\begin{quote}
Arun goes to the office by motorcycle.
\end{quote}

Feminine.

\begin{cousinweb}
\textbf{\hi{मोटरसाइकिल}} (\emph{moṭarsāikil}) is borrowed from English \emph{motorcycle}; like \hi{साइकिल} it is feminine in Hindi.
\end{cousinweb}

\begin{grammarlens}[title={What you've built}]
One more word for getting around.
\end{grammarlens}

\subsection*{Your turn}
\begin{itemize}
\raggedright
  \item \emph{Say it:} \emph{moṭarsāikil}
  \item \emph{Say it:} \emph{moṭarsāikil}, once more
  \item \emph{Say it:} \emph{Aruṇ moṭarsāikil se daftar jātā hai.} — Arun goes to the office by motorcycle.
  \item \emph{Recall:} say the Hindi for traffic, then the Hindi for a traffic jam, then say \emph{moṭarsāikil} again
\end{itemize}

\begin{tcolorbox}[breakable,colback=teal!4,colframe=teal!35!black,title={Before you move on}]
What does \hi{मोटरसाइकिल} mean? (\textquotedblleft{}A motorcycle\textquotedblright{}.) Say it once more.
\end{tcolorbox}
\section[\texorpdfstring{\hi{रेंगना} (reṅgnā)}{रेंगना (reṅgnā)}]{\hi{रेंगना} (reṅgnā) — to crawl}

\label{lesson:HI-C404-rengna}



\begin{quote}
Before the new one: say the Hindi for a traffic jam, then the Hindi for a motorcycle.
\end{quote}

\subsection*{You'll want to know: \hi{रेंगना}}
\textbf{\hi{रेंगना}} — \emph{reṅgnā} — \textquotedblleft{}to crawl\textquotedblright{}.

\begin{quote}
\hi{जाम} \hi{में} \hi{गाड़ियाँ} \hi{रेंग} \hi{रही} \hi{थीं।}
\end{quote}

\begin{quote}
\emph{jām meṁ gāṛiyā̃ reṅg rahī thīṁ.}
\end{quote}

\begin{quote}
The cars were crawling along in the jam.
\end{quote}

Past (it …): \hi{वह} \hi{साँप} \hi{घास} \hi{में} \hi{रेंगा}, \emph{vah sā̃p ghās meṁ reṅgā}.

\begin{cousinweb}
\textbf{\hi{रेंगना}} (\emph{reṅgnā}) is a tadbhava of Sanskrit \emph{riṅgati}, \textquotedblleft{}crawls\textquotedblright{} (CDIAL 10739), which lists Hindi \emph{rẽgnā}. A baby, a snake and a jammed lane of traffic all crawl.
\end{cousinweb}

\begin{grammarlens}[title={What you've built}]
One more word for getting around.
\end{grammarlens}

\subsection*{Your turn}
\begin{itemize}
\raggedright
  \item \emph{Say it:} \emph{reṅgnā}
  \item \emph{Say it:} \emph{reṅgnā}, once more
  \item \emph{Say it:} \emph{jām meṁ gāṛiyā̃ reṅg rahī thīṁ.} — The cars were crawling along in the jam.
  \item \emph{Recall:} say the Hindi for a traffic jam, then the Hindi for a motorcycle, then say \emph{reṅgnā} again
\end{itemize}

\begin{tcolorbox}[breakable,colback=teal!4,colframe=teal!35!black,title={Before you move on}]
What does \hi{रेंगना} mean? (\textquotedblleft{}To crawl\textquotedblright{}.) Say it once more.
\end{tcolorbox}
\section[\texorpdfstring{\hi{हाँकना} (hā̃knā)}{हाँकना (hā̃knā)}]{\hi{हाँकना} (hā̃knā) — to drive (animals); to drive (a cart)}

\label{lesson:HI-C404-hakna}



\begin{quote}
Before the new one: say the Hindi for a motorcycle, then the Hindi for to crawl.
\end{quote}

\subsection*{You'll want to know: \hi{हाँकना}}
\textbf{\hi{हाँकना}} — \emph{hā̃knā} — \textquotedblleft{}to drive (animals); to drive (a cart)\textquotedblright{}.

\begin{quote}
\hi{एक} \hi{किसान} \hi{सड़क} \hi{पर} \hi{भैंसें} \hi{हाँक} \hi{रहा} \hi{था।}
\end{quote}

\begin{quote}
\emph{ek kisān saṛak par bhaĩsẽ hā̃k rahā thā.}
\end{quote}

\begin{quote}
A farmer was driving buffaloes along the road.
\end{quote}

Past, with \hi{ने} (he or she …): \hi{उसने} \hi{गायें} \hi{हाँकीं}, \emph{usne gāyẽ hā̃kīṁ}.

\begin{cousinweb}
\textbf{\hi{हाँकना}} (\emph{hā̃knā}) is a tadbhava of Sanskrit \emph{hakkayati}, \textquotedblleft{}calls to\textquotedblright{} (CDIAL 13939), which lists Hindi \emph{hā̃knā}, \textquotedblleft{}to cry out, drive off\textquotedblright{}: cattle are driven by shouting at them.
\end{cousinweb}

\begin{grammarlens}[title={What you've built}]
One more word for getting around.
\end{grammarlens}

\subsection*{Your turn}
\begin{itemize}
\raggedright
  \item \emph{Say it:} \emph{hā̃knā}
  \item \emph{Say it:} \emph{hā̃knā}, once more
  \item \emph{Say it:} \emph{ek kisān saṛak par bhaĩsẽ hā̃k rahā thā.} — A farmer was driving buffaloes along the road.
  \item \emph{Recall:} say the Hindi for a motorcycle, then the Hindi for to crawl, then say \emph{hā̃knā} again
\end{itemize}

\begin{tcolorbox}[breakable,colback=teal!4,colframe=teal!35!black,title={Before you move on}]
What does \hi{हाँकना} mean? (\textquotedblleft{}To drive (animals); to drive (a cart)\textquotedblright{}.) Say it once more.
\end{tcolorbox}
